\chapter*{第三部分收束：从方法章走向案例项目}
\addcontentsline{toc}{chapter}{第三部分收束：从方法章走向案例项目}

Part III 已经给出了机器学习的基本方法，但方法章和真实案例之间还隔着一步：把“我会用模型”转化成“我能把模型嵌入一个科学问题”。这一步非常关键，因为真实案例很少按照教科书顺序出现。你拿到的往往是一份不完美的数据、一组含糊的标签、一个需要折中的评估指标，以及一个并不完全适合机器学习的问题。

\section*{一个最小模型记录}

进入案例项目之前，每次模型实验至少应留下下面几类记录：

\begin{itemize}
    \item \textbf{任务定义}：回归、分类、排序、聚类还是异常候选筛选；
    \item \textbf{数据入口}：使用了哪些文件、字段、派生特征和筛选条件；
    \item \textbf{baseline}：最简单可解释方法的结果；
    \item \textbf{模型与参数}：使用了什么模型，关键超参数是什么；
    \item \textbf{评估设计}：训练/验证/测试如何划分，指标为什么适合当前科学问题；
    \item \textbf{错误分析}：哪些样本错得最多，错误是否集中在某类物理条件或数据质量上；
    \item \textbf{解释边界}：结果能支持什么，不能支持什么。
\end{itemize}

这份记录不是行政负担，而是后续写报告、展示、答疑和复现实验的核心材料。

\section*{模型结果报告的五个固定问题}

延续 Part II 的 evidence packet 思路，一个机器学习结果也不应只报告一行分数。每次模型实验至少要回答五个固定问题：

\begin{enumerate}
    \item \textbf{任务是什么}：预测连续量、分类标签、发现结构还是筛选候选，科学目标是否明确；
    \item \textbf{数据入口是什么}：特征、标签、样本筛选、派生量和预处理是否来自可追溯流程；
    \item \textbf{baseline 是什么}：最简单可解释方法表现如何，复杂模型到底改进了什么；
    \item \textbf{评估设计是什么}：训练/验证/测试如何切分，指标是否匹配当前科学代价；
    \item \textbf{错误和边界是什么}：哪些样本失败，失败是否集中在某类物理条件、数据质量或外推区域。
\end{enumerate}

如果这五问答不出来，模型分数越高越危险，因为它可能只是把数据泄漏、选择效应或标签噪声包装成了“好结果”。如果这五问能答出来，即使模型很简单，它也已经是一个可以讨论、可以复查、可以进入案例项目的科学工作流。

\section*{这些记录怎样进入 Capstone}

Part III 的每份模型记录，到了 Part VI 会直接变成 capstone case template 里的几张卡片：

\begin{itemize}
    \item 任务定义会进入 Project card；
    \item 数据入口、筛选条件和预处理会进入 Data card 与 Reproducibility card；
    \item baseline、模型参数和训练设置会进入 Baseline card 与 Model card；
    \item 评估设计、错误分析和边界说明会进入 Evaluation card 与 Interpretation card；
    \item 数据泄漏检查、随机种子、版本记录和人工复核会进入 Integrity card。
\end{itemize}

因此，从这一部分开始，推荐把每一次实验都当作“未来 capstone 包的一张证据卡”来保存。不要等到期末才回忆模型怎么训练、数据怎么划分、为什么选择这个指标；那时很多关键判断已经很难复原。

\section*{从分数到样本}

Part III 反复强调一个习惯：不要只看平均分。平均分会把许多重要信息压扁，而科学问题常常正藏在被压扁的部分里。例如：

\begin{itemize}
    \item 回归模型的残差是否随颜色、亮度或红移系统变化；
    \item 分类模型是否总把稀有类别错分成常见类别；
    \item 聚类结果是否只是预处理尺度的产物；
    \item 异常检测候选是否来自坏点、缺失值或边界条件；
    \item 解释性图表是否揭示了模型依赖了不该依赖的特征。
\end{itemize}

因此，进入 Part IV 时，你应该带着一个朴素但强大的原则：\textbf{模型分数只是入口，样本级诊断才是科学解释的开始。}

\section*{Part III 最常见的误区}

回到案例之前，值得把方法章中最容易反复出现的错误集中列出来：

\begin{itemize}
    \item 先调模型，再回头定义任务；
    \item 没有 baseline，就直接比较复杂模型；
    \item 在全数据上做标准化、缺失值填补或特征选择，造成数据泄漏；
    \item 只报告准确率、MAE 或平均分，不检查错误样本；
    \item 把验证集反复用于选择模型，最后又把它当成无偏测试；
    \item 把特征重要性或局部解释误读为因果关系；
    \item 把无监督聚类结果直接命名成物理类别；
    \item 把异常检测候选直接当作新发现，而没有做数据质量和人工复核。
\end{itemize}

这些错误的共同点是：它们通常不会让代码报错，却会让科学解释变得不可靠。因此，Part III 的训练重点不是“避免程序崩溃”，而是避免生成看起来漂亮、实际站不住的结果。

\section*{给案例章节的阅读方式}

后面的案例章不需要每次都引入全新的算法。它们更像是把前面工具重新组合：

\begin{enumerate}
    \item Part II 提供数据理解和处理；
    \item Part III 提供模型流程和评估语言；
    \item Part IV 把两者放进真实天文/物理问题中；
    \item 每个案例都应产出可运行 notebook、结果图表、错误分析和有限结论。
\end{enumerate}

如果某个案例使用的模型很简单，这并不意味着它“不高级”。本科阶段真正值得训练的是：是否能把简单模型做对、解释清楚、知道边界，并为之后引入深度学习或现代 AI 留下可靠基线。
